
  \section{Statistical Evaluation Measures}\label{sec:stat-eval}
    The statistical metrics considered in this paper are used to
    evaluate either the realism of synthetic market data to its
    trained history or generative novelty in market scenarios.
    The following subsections describes the metrics considered and
    the evaluatiom framework applied.
    
    \subsection{Statistical Measures}\label{subsec:stats-measure}
    As described in \ref{limitation:metrics}, our contribution 
    \ref{contribute:1} is based on a well-established taxonomy framework
    \cite{Stenger24} while introducing our industry-specific stylized
    facts to holistically evaluate the characteristics of synthetic
    market data. We divide our statistical metrics
    \subsubsection{Fidelity Measures}
    These metrics quantify how well synthetic data captures key 
    marginal and second-order properties of real data. We employ
    the following feature-based metrics:

      \begin{enumerate}[label=\textbf{[M\arabic*]}]
        \item \label{metric:mdd} \textbf{Marginal Distribution 
        Difference (MDD):}  
        To measure the distance between empirical marginal distributions
        of real versus synthetic data, we use the 1-Wasserstein distance
        (a.k.a. Earth Mover's distance). This was chosen over the 
        popular Kullback-Leibler divergence or the derivatives
        Jensen-Shannon divergence due to its better numerical stability
        and ability to capture distributional differences even when
        supports do not overlap \cite{Arjovsky17}.
        \[
        \begin{aligned}
          \mathrm{MDD} &= W_1\big(\widehat{F}_{\text{real}}, 
          \widehat{F}_{\text{synth}}\big), \\
          W_1(F, G) &= \int_0^1 \big|F^{-1}(u) - G^{-1}(u)\big|\,du,
          \end{aligned}
          \]
          where \(\widehat{F}_{\text{real}}\) and \(\widehat{F}_{\text{synth}}\) 
          are empirical CDFs.
          
          \end{enumerate}
          The following four metrics capture differences in the first four
          statistical (central) moments between real and synthetic data.
          \begin{enumerate}[label=\textbf{[M\arabic*]}, resume]
          
          \item \label{metric:md} \textbf{Mean Difference (MD):}  
          \[
          \mathrm{MD} = \big|\mu_{\text{real}} - \mu_{\text{synth}}\big|.
          \]
          
          \item \label{metric:sdd} \textbf{Standard Deviation Difference (SDD):}  
          \[
          \mathrm{SDD} = \big|\sigma_{\text{real}} - \sigma_{\text{synth}}\big|.
          \]

          \item \label{metric:skewness} \textbf{Skewness Difference (SD):}  
          \begin{align*}
          \mathrm{SD} =& \big|\gamma(r_{\text{real}}) - \gamma(r_{\text{synth}})\big|, \\
          \gamma(r) =& \frac{\mathbb{E}[(r - \mu)^3]}{(\sigma)^3}.
          \end{align*}
          
          \item \label{metric:kd} \textbf{Kurtosis Difference (KD):}  
          \begin{align*}
          \mathrm{KD} &= \big|\kappa^{\mathrm{ex}}(r_{\text{real}}) 
          - \kappa^{\mathrm{ex}}(r_{\text{synth}})\big|, \\
          \kappa^{\mathrm{ex}}(r) &= \frac{\mathbb{E}[(r - \mu)^4]}{(\sigma)^4} - 3.
          \end{align*}
          
          \end{enumerate}
          
          \subsubsection{Visualization Methods} \label{sec:visualization_methods}
          Visual tools complement numeric metrics for face-validity and 
          communication \cite{Ang23}.
          \begin{enumerate}[label=\textbf{[M\arabic*]}, resume]
          \item \label{metric:tsne} \textbf{t-SNE:}  
          2D embeddings of window-level features for real 
          versus synthetic overlap and cluster structure.
          
          \item \label{metric:dist} \textbf{Distribution:}  
          Histograms of returns to visualize empirical 
          distribution.
          \end{enumerate}
          
          \subsubsection{Diversity Measures}
          We are also interested if the SDGFTS can generate data
          not observed but follow the same distribution as the real data.
          Therefore, we include the following intra-class distance
          metrics. The inclusion of
          these metrics are chose to ensure that mode collapse 
          is prevented. 

          
          \begin{enumerate}[label=\textbf{[M\arabic*]}, resume]
          \item \label{metric:ed} \textbf{Intra-Class Euclidean 
          Distance (ICD-ED):}  
          \[
          \mathrm{ICD\text{-}ED} = \frac{2}{R(R-1)} 
          \sum_{1 \le a < b \le R} \left\| \textbf{r}_a - 
          \textbf{r}_b \right\|_2,
          \]
          where \(\textbf{r}_a \in \mathbb{R}^{L}\) is the \(a\)-th 
          sample.
          
          \item \label{metric:dtw} \textbf{Intra-Class Dynamic Time 
          Warping (ICD-DTW):}  
          \[
          \mathrm{ICD\text{-}DTW} = \frac{2}{R(R-1)} 
          \sum_{1 \le a < b \le R} d_{\mathrm{DTW}}\bigl(
          \textbf{r}_a, \textbf{r}_b\bigr).
          \]
          \end{enumerate}
          
          \subsubsection{Efficiency Measures}
          Generation efficiency is crucial for practical applications,
          especially in high-frequency trading or real-time risk
          management, where rapid data synthesis is required.
          Thus, we include the following efficiency metric \cite{Stenger24}.
          \begin{enumerate}[label=\textbf{[M\arabic*]}, resume]
          \item \label{metric:gentime} \textbf{Generation Time:}  
          Wall-clock time to generate \(S\) samples:
          \[
          T_{\mathrm{gen}} = t_1 - t_0 \quad (\text{seconds for } 1,000
          \text{ samples}).
          \]
          \end{enumerate}
          
          \subsubsection{Stylized Facts Measures} \label{sec:stylized-facts}
          Canonical stylized facts:
          
          \begin{enumerate}[label=\textbf{[M\arabic*]}, resume]

          \item \label{metric:acd} \textbf{Autocorrelation Difference 
          (ACD):}  
          Absolute difference in lag-1 autocorrelation:
          \[
          \begin{aligned}
          \rho(1) &= \frac{\sum_{t=2}^{L} (r_t - \bar{r})
          (r_{t-1} - \bar{r})}
               {\sum_{t=1}^{L} (r_t - \bar{r})^2}, \\
          \mathrm{ACD} &= \big|\rho_{\text{real}}(1) - 
          \rho_{\text{synth}}(1)\big|.
          \end{aligned}
          \]
          
          \item \label{metric:volclustering} \textbf{Volatility 
          Clustering (VC):}  
          Lag-1 autocorrelation of squared or absolute returns measures 
          volatility clustering, a stylized fact where large price movements 
          tend to cluster in time \cite{Cont01}:
          \[
          \begin{aligned}
          \rho_{r^2}(1) &= \mathrm{Corr}\big((r_t)^2, 
          (r_{t-1})^2\big), \\
          \rho_{|r|}(1) &= \mathrm{Corr}\big(|r_t|, 
          |r_{t-1}|\big).
          \end{aligned}
          \]

        \item \label{metric:slow_decay} \textbf{Long Memory / Slow Decay (LMSD):}  
        To quantify long-range dependence in volatility, we follow
        Cont (2001) \cite{Cont01} and study the autocorrelation of
        absolute (or squared) returns:
        \[
        C_\alpha(\tau) \,=\, \mathrm{corr}\!\big(|r_{t+\tau}|^\alpha,\,
        |r_t|^\alpha\big), \qquad \alpha \in \{1,2\}.
        \]
        Empirically, this correlation decays slowly according to a
        power law,
        \[
        C_\alpha(\tau) \,\sim\, \frac{A}{\tau^{\beta}}, \qquad
        \beta \in [0.2,\,0.4],
        \]
        indicating persistent volatility fluctuations (long memory).
        For evaluation, we estimate the decay exponent by a
        log--log linear fit over lags 
        \(\tau \in \{1,\ldots, \tau_{\max}\}\) where
        \(C_\alpha(\tau) > 0\):
        \[
        \widehat{\beta} \,=\, -\,\text{slope}\Big(\log C_\alpha(\tau)
        \;\text{vs.}\; \log \tau\Big).
        \]
        Our metric is the absolute difference between decay exponents
        estimated on real and synthetic data:
        \[
        \mathrm{LMSD} \,=\, \big|\widehat{\beta}_{\text{real}} -
        \widehat{\beta}_{\text{synth}}\big|.
        \]
        Unless otherwise stated, we set \(\alpha=1\) (absolute returns)
        and report \(\widehat{\beta}\) together with the goodness-of-fit
        of the linear regression.
        \end{enumerate}
    
    \subsection{Statistical Evaluation Framework}
    We compare the statistical measures outlined above between the synthetic
    market data to the trained historical dataset. Consider a market generator
    $G_j, j\in\mathcal{J}$ trained on the dataset $D^{train}$ and its generated data
    $\mathcal{D}_{G_j}$ as specified in \ref{subsec:generator-training},
    we can apply an evaluation metric E, as mentioned in \ref{subsec:stats-measure}, defined an measurable function 
    $\mathbb{R}^{T\times K}\to \mathbb{R}$.